\section{Abstract}

Agriculture has been the cornerstone of India’s economy and culture for centuries. Often referred to as the backbone of rural India, it not only sustains a large section of the population but also contributes significantly to the country’s GDP and employment. Despite rapid industrialization and the growth of the service sector, agriculture continues to play a crucial role in the socio-economic fabric of the country.

India, with its diverse agro-climatic zones and rich agricultural heritage, is home to a wide range of food and non-food crops. The country’s agricultural landscape is shaped by varying soil types, climatic conditions, and traditional cultivation practices, resulting in the growth of crops that sustain millions of livelihoods. 

In this article our main goal will be :
\vspace{2 mm}

\hspace{-5 mm}\textbf{i)} {\color{blue}To identify the states of India which has largest amount of production of different crops (rice,wheat,jute etc.) }

\vspace{1 mm}
\hspace{-5 mm}\textbf{ii)} {\color{blue}explore any potential correlations between crop yield and different climatic factors(like rainfall , temperature etc)}

\vspace{1 mm}
\hspace{-5 mm}\textbf{iii)} {\color{blue} To analyze overall trends over time and to predict this year possible crop yield after concerning different factors.}

\vspace{1 mm}

\hspace{-5 mm}\textbf{iv)}{\color{blue} We are also interested to examine different types of problems in Indian agricultural system and come up with ways to reduce these problems via policy measures or farmers outreach programs.}
